\begin{table}[htbp]
\centering
\caption{Dez municípios com maior predominância masculina e feminina, Sergipe, 2022}
\label{tab:razao-sexo-top10}
\small
\begin{tabular}{r l l r}
\toprule
Posição & Município & Predominância & Homens por 100 mulheres \\
\midrule
\multicolumn{4}{l}{\textbf{Maior predominância masculina}} \\
1 & Canhoba & Masculina & 106,71 \\
2 & Gararu & Masculina & 106,17 \\
3 & Itabi & Masculina & 102,95 \\
4 & Brejo Grande & Masculina & 102,77 \\
5 & Poço Redondo & Masculina & 102,32 \\
6 & Porto da Folha & Masculina & 101,90 \\
7 & Areia Branca & Masculina & 101,82 \\
8 & Japoatã & Masculina & 101,70 \\
9 & Santa Rosa de Lima & Masculina & 101,59 \\
10 & Gracho Cardoso & Masculina & 101,38 \\
\midrule
\multicolumn{4}{l}{\textbf{Maior predominância feminina}} \\
1 & Aracaju & Feminina & 84,81 \\
2 & Rosário do Catete & Feminina & 87,89 \\
3 & General Maynard & Feminina & 89,22 \\
4 & Nossa Senhora do Socorro & Feminina & 90,11 \\
5 & Carmópolis & Feminina & 91,13 \\
6 & Estância & Feminina & 91,48 \\
7 & Umbaúba & Feminina & 91,80 \\
8 & Barra dos Coqueiros & Feminina & 91,83 \\
9 & Laranjeiras & Feminina & 91,95 \\
10 & Propriá & Feminina & 91,97 \\
\bottomrule
\end{tabular}
\fonte{Elaboração própria com base nos microdados do Censo Demográfico 2022.}
\nota{Razão de sexo calculada como o número de homens para cada 100 mulheres.}
\end{table}
